\section{Methods}

\subsection{Self-describing decoding}

The RS(255,127)-protected header records the format and method identifiers, ECC layout, image dimensions, layer codeword counts, payload length, interleaver and scrambler identifiers, payload CRC32, and a SHA-256 digest of the canonical configuration. Base and Detail use separate deterministic randomisation streams derived from the master seed, pair, method, layer, and purpose.

The receiver validates the header before accepting either layer. A secret image is produced only when the required Reed--Solomon decodes succeed and the complete payload CRC matches. When decoding fails, the record retains the stage, layer, codeword index, and corrected-symbol counts. This makes partial correction observable without presenting uncertain pixels as recovered data.

\subsection{Audited transform interpretation}

The final profile uses 9--7 pyramid filters, PKVA directional filtering, four multiscale levels, and independent P3+P4 range coordinates. The adapter verifies that required functions resolve inside the locked PDFB toolbox and records the toolbox, bridge, band inventory, and transform fingerprint. Reconstruction, capacity, writability, self-gain, cross-talk, and off-target-energy gates are reported in Supplementary Note S2. The profile is an explicit independent interpretation rather than an inferred copy of the source authors' undisclosed configuration.

\subsection{Factorial design and channels}

The experiment separates adaptive allocation A from unequal protection D. Figure~\ref{fig:factorial} shows the four controlled configurations.

\begin{figure}[t]
\centering
\begin{tikzpicture}[font=\small,
  cell/.style={draw, rounded corners=2pt, minimum width=39mm, minimum height=16mm, align=center}]
\node[cell, fill=palegray] (c0) at (0,1.1) {\textbf{C0}\\Fixed allocation\\Symmetric coding};
\node[cell, fill=accentorange!10] (c2) at (4.4,1.1) {\textbf{C2}\\Fixed allocation\\Unequal protection};
\node[cell, fill=baseblue!10] (c1) at (0,-1.1) {\textbf{C1}\\Adaptive allocation\\Symmetric coding};
\node[cell, fill=softgreen!12] (c3) at (4.4,-1.1) {\textbf{C3}\\Adaptive allocation\\Unequal protection};
\node[font=\bfseries] at (2.2,2.35) {Unequal protection D};
\node at (0,1.95) {off};
\node at (4.4,1.95) {on};
\node[font=\bfseries, rotate=90] at (-2.65,0) {Adaptive allocation A};
\node at (-2.05,1.1) {off};
\node at (-2.05,-1.1) {on};
\end{tikzpicture}
\caption{The \(2\times2\) design isolates adaptive allocation and unequal protection before testing their combination in C3.}
\label{fig:factorial}
\end{figure}

Four fixed traceability pairs yield 16 embeddings. Mandatory evaluation includes clean output, JPEG quality 70, Gaussian variance 10, and salt-and-pepper density 0.03. Predeclared triggers activate hard checks at JPEG quality 50, Gaussian variance 15, and salt-and-pepper density 0.05, producing 64 mandatory and 24 conditional rows. The pair, payload, transform fingerprint, PSNR rule, and channel realisation remain aligned across methods.

Recorded outcomes include cover--stego quality, raw BER, Base and Detail BER, header and payload validity, valid-secret quality, selected strength, runtime, memory, and immutable object identifiers. The primary validity measure is effective unrecovered bit rate, \(\eur\): zero denotes valid complete recovery; one denotes failure of the complete recovery contract.

\subsection{Execution and scope}

Each embedding and channel evaluation is stored as a content-addressed object and reused only after schema and SHA-256 validation. A runtime gate terminates a real worker with SIGKILL, verifies that completed objects remain unchanged, resumes unfinished work, and validates the exported archive. The study is case-bounded: four fixed pairs support descriptive contrasts, not population-level claims. Extraction is semi-blind, geometric attacks and steganalysis are outside the final factorial evidence, and deterministic scrambling is not treated as encryption. Supplementary Notes S4--S6 provide the statistical, runtime, and claim-boundary details.
